% !TeX root = ../../main.tex
\section{Legge della Probabilità Totale}
Il teorema permette di calcolare la probabilità di un evento $A$ sfruttando una partizione dello spazio campionario $\Omega$.

\subsection{Ipotesi: La Partizione}

Sia $\{E_i\}_{i=1}^M$ una famiglia di insiemi che costituisce una \textbf{partizione} di $\Omega$. Affinché sia una partizione, devono valere due condizioni fondamentali:

\begin{enumerate}
    \item \textbf{Esaustività:} L'unione degli insiemi copre l'intero spazio campionario.
    \[ \bigcup_{i=1}^M E_i = \Omega \]
    
    \item \textbf{Disgiunzione a coppie:} Gli insiemi non hanno elementi in comune tra loro.
    \[ E_i \cap E_j = \emptyset, \quad \forall i \neq j \]
\end{enumerate}

\subsection{Scomposizione dell'evento A}

Consideriamo un evento generico $A \subseteq \Omega$.
\\Possiamo riscrivere $A$ come l'intersezione di se stesso con l'intero spazio campionario $\Omega$ (poiché intersecare con l'insieme universo non cambia l'insieme):
\[ A = A \cap \Omega \]

Sostituendo $\Omega$ con la definizione di partizione (punto 1), otteniamo:
\[ A = A \cap \left( \bigcup_{i=1}^M E_i \right) \]

Applicando la \textbf{proprietà distributiva} dell'intersezione rispetto all'unione, portiamo l'intersezione dentro la sommatoria insiemistica:
\[ A = \bigcup_{i=1}^M (A \cap E_i) \]

\paragraph{Nota fondamentale}
Poiché gli insiemi $E_i$ sono disgiunti tra loro, anche le loro intersezioni con $A$ (ovvero i sottoinsiemi $A \cap E_i$) sono disgiunte a coppie.
\[ (A \cap E_i) \cap (A \cap E_j) = \emptyset, \quad \forall i \neq j \]

\subsection{Passaggio alla Probabilità (Assioma di Additività)}

Calcoliamo ora la probabilità di $A$, ovvero $\mathbb{P}(A)$.
\\Sfruttando la scomposizione appena trovata:
\[ \mathbb{P}(A) = \mathbb{P}\left( \bigcup_{i=1}^M (A \cap E_i) \right) \]

Per il \textbf{terzo assioma di Kolmogorov} (additività finita), poiché gli eventi $(A \cap E_i)$ sono disgiunti, la probabilità dell'unione è uguale alla somma delle probabilità:
\[ \mathbb{P}(A) = \sum_{i=1}^M \mathbb{P}(A \cap E_i) \]

\subsection{Applicazione della Probabilità Condizionata}

Dalla definizione di probabilità condizionata sappiamo che:
\[ \mathbb{P}(A \mid E_i) = \frac{\mathbb{P}(A \cap E_i)}{\mathbb{P}(E_i)} \quad (\text{con } \mathbb{P}(E_i) > 0) \]

Invertendo la formula, possiamo esprimere la probabilità congiunta come prodotto:
\[ \mathbb{P}(A \cap E_i) = \mathbb{P}(A \mid E_i) \cdot \mathbb{P}(E_i) \]

Sostituendo quest'ultima espressione nella sommatoria del punto 3, otteniamo la formula della Legge della Probabilità Totale:

\begin{equation}
\boxed{\mathbb{P}(A) = \sum_{i=1}^M \mathbb{P}(A \mid E_i)\mathbb{P}(E_i)}
\end{equation}

\paragraph{Interpretazione analitica}
Il teorema afferma che la probabilità marginale di $A$ è una \textbf{media ponderata} delle probabilità condizionate $\mathbb{P}(A \mid E_i)$, dove i pesi sono dati dalle probabilità degli eventi della partizione $\mathbb{P}(E_i)$.

